\documentclass[11pt,a4paper]{article}
% preamble.tex — shared setup for human-verification documents
% Compile each document with:  tectonic <File>.tex

\usepackage[T1]{fontenc}
\usepackage[utf8]{inputenc}
\usepackage{lmodern}
\usepackage{amssymb,amsmath,amsthm}
\usepackage{mathtools}
\usepackage[dvipsnames]{xcolor}
\usepackage{listings}
\usepackage{hyperref}
\usepackage[margin=2.5cm]{geometry}
\usepackage{parskip}

% ── Colours ──────────────────────────────────────────────────────────
\definecolor{leanBg}{HTML}{F7F7F0}
\definecolor{haskellBg}{HTML}{F0F4F7}
\definecolor{keyword}{HTML}{7B3294}
\definecolor{comment}{HTML}{5A8C32}
\definecolor{noteFg}{HTML}{B35900}

% ── Lean 4 listings language ─────────────────────────────────────────
% NOTE: Lean listing blocks use ASCII approximations of unicode symbols
% (e.g. -> for arrow, o. for compose).  The literate key substitutes
% these with proper math-mode renderings in the PDF.
\lstdefinelanguage{Lean4}{
  morekeywords={def,theorem,structure,where,namespace,end,import,fun,
                let,have,show,by,exact,intro,cases,rfl,subst,obtain,
                funext,simp,rw,apply,variable,Type,Prop,instance,
                section,open,class,extends,deriving,sorry,noncomputable},
  sensitive=true,
  morecomment=[l]{--},
  morecomment=[n]{/-}{-/},
  morestring=[b]",
}

% ── Lean listing style ───────────────────────────────────────────────
\lstdefinestyle{lean}{
  language=Lean4,
  backgroundcolor=\color{leanBg},
  basicstyle=\small\ttfamily,
  keywordstyle=\color{keyword}\bfseries,
  commentstyle=\color{comment}\itshape,
  breaklines=true,
  frame=single,
  framerule=0.4pt,
  xleftmargin=4pt,
  xrightmargin=4pt,
  aboveskip=6pt,
  belowskip=4pt,
  literate=
    {->}{{$\to$}}2
    {=>}{{$\Rightarrow$}}2
    {<|}{{$\langle$}}1
    {|>}{{$\rangle$}}1
    {==>}{{$\Longrightarrow$}}3
    { o. }{{$\;\circ\;$}}3
    {forall}{{$\forall$}}1
    {:=}{{$\,:=\,$}}2
    {!=}{{$\neq$}}2
    {<=}{{$\leq$}}2
    {>=}{{$\geq$}}2,
}

% ── Haskell listing style ───────────────────────────────────────────
\lstdefinestyle{haskell}{
  language=Haskell,
  backgroundcolor=\color{haskellBg},
  basicstyle=\small\ttfamily,
  keywordstyle=\color{keyword}\bfseries,
  commentstyle=\color{comment}\itshape,
  breaklines=true,
  frame=single,
  framerule=0.4pt,
  xleftmargin=4pt,
  xrightmargin=4pt,
  aboveskip=6pt,
  belowskip=4pt,
  literate=
    {->}{{$\to$}}2
    {=>}{{$\Rightarrow$}}2,
}

% ── Theorem environments ────────────────────────────────────────────
\theoremstyle{definition}
\newtheorem{defn}{Definition}[section]
\newtheorem{thm}[defn]{Theorem}
\newtheorem{prop}[defn]{Proposition}

\theoremstyle{remark}
\newtheorem*{fidelitynote}{Fidelity note}

% ── Convenience macros ───────────────────────────────────────────────
\newcommand{\lean}[1]{\texttt{#1}}      % inline Lean identifier
\newcommand{\hask}[1]{\texttt{#1}}      % inline Haskell identifier
\newcommand{\Ext}{\mathrm{Ext}}
\newcommand{\Pos}{\mathrm{Pos}}
\newcommand{\Shape}{\mathrm{Shape}}

\usepackage{stmaryrd}

\title{Human Verification: \texttt{Sequential}\\[4pt]
  \large The Sequential Operator $\triangleleft$ and its Monoidal Structure}
\author{Machine-generated from \texttt{lean/Sequential.lean} for review by Neil Ghani}
\date{July 2026}

\begin{document}
\maketitle

\begin{abstract}
This document presents every definition and theorem \emph{statement} from the
Lean~4 module \lean{Sequential} in three parallel forms: the original Lean,
a readable Haskell rendering, and precise mathematical English.  Proofs are
omitted --- they are machine-checked by Lean.

The module is a near-duplicate of \lean{Composition.lean}, re-exposing the
same monoidal structure under the name ``sequential operator'' rather than
``composition product''.  The mathematical content is identical; only the
terminology and the names of a few definitions differ (\lean{Container.seq}
vs \lean{Container.comp}, \lean{seqExt} vs \lean{compExt}).  A fidelity note
at the end records the relationship.
\end{abstract}

\tableofcontents
\bigskip

%===================================================================
\section{Overview}
%===================================================================

The sequential operator $\triangleleft$ on containers has the same formula as
the composition product of \lean{Composition.lean}: for $F = S \triangleleft
P$ and $G = T \triangleleft Q$,
\begin{itemize}
  \item shapes of $G \triangleleft F$: $\Sigma_{t : T}(Q(t) \to S)$, and
  \item positions at $(t, f)$: $\Sigma_{q : Q(t)} P(f(q))$.
\end{itemize}
The extension satisfies $\llbracket G \triangleleft F \rrbracket = \llbracket
G \rrbracket \circ \llbracket F \rrbracket$.  The monoidal unit is the same
identity container $I$.  The coherence isomorphisms (left/right unitors,
associator) have identical statements and proofs.

The module's docstring frames the same structure as the ``sequential
operator'' to emphasise the operational reading: $G \triangleleft F$ is ``run
$F$, then run $G$ on the result'', and a $\triangleleft$-comonoid is a
directed container.

%===================================================================
\section{The sequential operator}
%===================================================================

\subsection{Definition: \lean{Container.seq}}

\begin{defn}[Sequential operator]
For containers $G$ and $F$, their \textbf{sequential composition}
$G \triangleleft F$ (``$G$ outside, $F$ inside'') is:
\[
  \mathrm{Shape}(G \triangleleft F) = \Sigma_{t : G.\mathrm{Shape}}
    \bigl(G.\mathrm{Pos}(t) \to F.\mathrm{Shape}\bigr),
\]
\[
  \mathrm{Pos}(G \triangleleft F)(t, f) = \Sigma_{q : G.\mathrm{Pos}(t)}
    F.\mathrm{Pos}(f(q)).
\]
The extension satisfies $\llbracket G \triangleleft F \rrbracket X \cong
\llbracket G \rrbracket(\llbracket F \rrbracket X)$.
\end{defn}

\paragraph{Lean.}
\begin{lstlisting}[style=lean]
def Container.seq (G F : Container) : Container where
  Shape := (t : G.Shape) * (G.Pos t -> F.Shape)
  Pos := fun tf => (q : G.Pos tf.1) * F.Pos (tf.2 q)

scoped infixr:70 " seq " => Container.seq
-- (written with the unicode diamond operator in Lean source)
\end{lstlisting}

\paragraph{Haskell.}
\begin{lstlisting}[style=haskell]
-- G `seq` F: shapes = (t, p_g t -> s_f), positions = (p_g t, p_f (f q))
-- Same as Container.comp; see fidelity note.
data Seq g p_g f p_f x = Seq g (p_g -> f) (p_g -> p_f -> x)
\end{lstlisting}

\begin{fidelitynote}
\lean{Container.seq} and \lean{Container.comp} (from \lean{Composition.lean})
are definitionally identical --- same formula, same type.  The duplication
appears to be a development artifact, with \lean{Sequential.lean} providing
an alternative narrative framing of the same monoidal structure.  The
verifier should treat the two files as presenting the same mathematics.
\end{fidelitynote}


\subsection{Definition: \lean{Container.I}}

\begin{defn}[Identity container]
The identity container $I$ has one shape and one position; its extension is
the identity functor.
\end{defn}

\paragraph{Lean.}
\begin{lstlisting}[style=lean]
def Container.I : Container where
  Shape := Unit
  Pos := fun _ => Unit
\end{lstlisting}

\paragraph{Haskell.}
\begin{lstlisting}[style=haskell]
newtype Id x = Id x   -- one shape, one position
\end{lstlisting}

%===================================================================
\section{Extension isomorphism}
%===================================================================

\subsection{Definitions: \lean{Container.seqExt} and \lean{Container.seqExtInv}}

\begin{defn}[Extension isomorphism for the sequential operator]
There is a natural bijection
\[
  \mathrm{seqExt} : \llbracket G \triangleleft F \rrbracket X \;\xrightarrow{\;\sim\;}
  \llbracket G \rrbracket\bigl(\llbracket F \rrbracket X\bigr),
\]
\[
  \mathrm{seqExtInv} : \llbracket G \rrbracket\bigl(\llbracket F \rrbracket X\bigr)
  \;\xrightarrow{\;\sim\;} \llbracket G \triangleleft F \rrbracket X.
\]
The formulas are identical to those of \lean{compExt} / \lean{compExtInv}.
\end{defn}

\paragraph{Lean.}
\begin{lstlisting}[style=lean]
def Container.seqExt {G F : Container} {X : Type} :
    Ext (G seq F) X -> Ext G (Ext F X) :=
  fun p => <|p.1.1, fun q => <|p.1.2 q, fun r => p.2 <|q, r|>|>|>

def Container.seqExtInv {G F : Container} {X : Type} :
    Ext G (Ext F X) -> Ext (G seq F) X :=
  fun p => <|<|p.1, fun q => (p.2 q).1|>, fun qr => (p.2 qr.1).2 qr.2|>
\end{lstlisting}

\paragraph{Haskell.}
\begin{lstlisting}[style=haskell]
seqExt :: Seq g p_g f p_f x -> (g, p_g -> (f, p_f -> x))
seqExt (Seq t assign lab) = (t, \q -> (assign q, \r -> lab q r))

seqExtInv :: (g, p_g -> (f, p_f -> x)) -> Seq g p_g f p_f x
seqExtInv (t, m) = Seq t (fst . m) (\q r -> snd (m q) r)
\end{lstlisting}


\subsection{Theorems: \lean{seqExtInv\_seqExt} and \lean{seqExt\_seqExtInv}}

\begin{thm}[Extension isomorphism is an isomorphism]
\[
  \mathrm{seqExtInv} \circ \mathrm{seqExt} = \mathrm{id},
  \qquad
  \mathrm{seqExt} \circ \mathrm{seqExtInv} = \mathrm{id}.
\]
Both hold definitionally (by \lean{rfl} after \lean{cases}).
\end{thm}

\paragraph{Lean.}
\begin{lstlisting}[style=lean]
theorem Container.seqExtInv_seqExt {G F : Container} {X : Type}
    (p : Ext (G seq F) X) :
    Container.seqExtInv (Container.seqExt p) = p

theorem Container.seqExt_seqExtInv {G F : Container} {X : Type}
    (p : Ext G (Ext F X)) :
    Container.seqExt (Container.seqExtInv p) = p
\end{lstlisting}

\paragraph{Haskell.}
\begin{lstlisting}[style=haskell]
-- seqExtInv . seqExt = id   (and vice versa)
-- Both hold by eta-laws for pairs.
\end{lstlisting}

%===================================================================
\section{Container isomorphisms}
%===================================================================

\subsection{Structure: \lean{ContainerIso}}

\begin{defn}[Container isomorphism]
An isomorphism in $\mathbf{Cont}$: a pair of mutually inverse container
morphisms.  (Identical definition to \lean{Composition.lean}.)
\end{defn}

\paragraph{Lean.}
\begin{lstlisting}[style=lean]
structure ContainerIso (C D : Container) where
  hom : ContainerMorphism C D
  inv : ContainerMorphism D C
  hom_inv : hom.comp inv = ContainerMorphism.id C
  inv_hom : inv.comp hom = ContainerMorphism.id D
\end{lstlisting}

\paragraph{Haskell.}
\begin{lstlisting}[style=haskell]
data ContainerIso c d = ContainerIso
  { hom :: Morph c d
  , inv :: Morph d c
  -- hom_inv and inv_hom are equational laws
  }
\end{lstlisting}

%===================================================================
\section{Monoidal coherences}
%===================================================================

\subsection{Definition: \lean{Container.leftUnitor}}

\begin{defn}[Left unitor]
$\lambda_F : I \triangleleft F \cong F$.  A shape $(\langle\rangle, f)$ of
$I \triangleleft F$ bijects with $F.\mathrm{Shape}$ via $f \mapsto f(\langle\rangle)$.
\end{defn}

\paragraph{Lean.}
\begin{lstlisting}[style=lean]
def Container.leftUnitor (F : Container) :
    ContainerIso (Container.I seq F) F where
  hom :=
    { onShapes := fun tf => tf.2 ()
      onPos := fun _ p => <|(), p|> }
  inv :=
    { onShapes := fun s => <|(), fun _ => s|>
      onPos := fun _ qp => qp.2 }
  hom_inv := rfl
  inv_hom := rfl
\end{lstlisting}

\paragraph{Haskell.}
\begin{lstlisting}[style=haskell]
leftUnitor :: ContainerIso (I `Seq` F) F
-- Same as in Composition.lean.
\end{lstlisting}


\subsection{Definition: \lean{Container.rightUnitor}}

\begin{defn}[Right unitor]
$\rho_F : F \triangleleft I \cong F$.  Shapes project onto $F.\mathrm{Shape}$;
positions collapse $\Sigma_{p : F.\mathrm{Pos}(s)} \mathrm{Unit} \cong F.\mathrm{Pos}(s)$.
\end{defn}

\paragraph{Lean.}
\begin{lstlisting}[style=lean]
def Container.rightUnitor (F : Container) :
    ContainerIso (F seq Container.I) F where
  hom :=
    { onShapes := fun sg => sg.1
      onPos := fun _ p => <|p, ()|> }
  inv :=
    { onShapes := fun s => <|s, fun _ => ()|>
      onPos := fun _ pu => pu.1 }
  hom_inv := rfl
  inv_hom := rfl
\end{lstlisting}

\paragraph{Haskell.}
\begin{lstlisting}[style=haskell]
rightUnitor :: ContainerIso (F `Seq` I) F
-- Same as in Composition.lean.
\end{lstlisting}


\subsection{Definition: \lean{Container.associator}}

\begin{defn}[Associator]
$\alpha_{H,G,F} : (H \triangleleft G) \triangleleft F \cong H \triangleleft (G \triangleleft F)$.
On shapes: curries the inner assignment.  On positions: re-associates the
nested dependent sum $((u, v), p) \leftrightarrow (u, (v, p))$.  No transport
needed (position fibres agree definitionally).
\end{defn}

\paragraph{Lean.}
\begin{lstlisting}[style=lean]
def Container.associator (H G F : Container) :
    ContainerIso ((H seq G) seq F) (H seq (G seq F)) where
  hom :=
    { onShapes := fun xk =>
        <|xk.1.1, fun u => <|xk.1.2 u, fun v => xk.2 <|u, v|>|>|>
      onPos := fun _ uvp => <|<|uvp.1, uvp.2.1|>, uvp.2.2|> }
  inv :=
    { onShapes := fun am =>
        <|<|am.1, fun u => (am.2 u).1|>, fun y => (am.2 y.1).2 y.2|>
      onPos := fun _ yp => <|yp.1.1, yp.1.2, yp.2|> }
  hom_inv := ...   -- ContainerMorphism.ext' rfl (cases + rfl)
  inv_hom := ...   -- ContainerMorphism.ext' rfl (cases + rfl)
\end{lstlisting}

\paragraph{Haskell.}
\begin{lstlisting}[style=haskell]
associator :: ContainerIso ((H `Seq` G) `Seq` F) (H `Seq` (G `Seq` F))
-- Same formulas as in Composition.lean.
\end{lstlisting}

%===================================================================
\section{The compositional payoff}
%===================================================================

A \textbf{$\triangleleft$-comonoid} in $(\mathbf{Cont}, \triangleleft, I)$ is
a container $C$ with a counit $\varepsilon : C \to I$ and a comultiplication
$\delta : C \to C \triangleleft C$, satisfying the comonoid laws.  This is
exactly a \textbf{directed container} (a comonad on $\mathbf{Set}$ whose
functor is a container).  The sequential framing makes the operational
reading explicit: a directed container is a container that can be
``stepped'', with the comonad structure recording how steps compose.

%===================================================================
\section{Relationship to \lean{Composition.lean}}
%===================================================================

\begin{fidelitynote}
\lean{Sequential.lean} and \lean{Composition.lean} define mathematically
identical structures: \lean{Container.seq} $=$ \lean{Container.comp},
\lean{seqExt} $=$ \lean{compExt}, and so on.  The files differ only in
terminology and narrative framing:
\begin{itemize}
  \item \lean{Composition.lean} emphasises $\llbracket G \triangleleft F
    \rrbracket = \llbracket G \rrbracket \circ \llbracket F \rrbracket$ as
    functor composition;
  \item \lean{Sequential.lean} emphasises the sequential reading and the
    comonoid/directed-container connection.
\end{itemize}
In a production codebase these would be merged into one file.  As written,
they coexist by declaring the notation in separate scopes; the notation
\texttt{infixr:70 " $\triangleleft$ "} is declared twice, once per file.
\end{fidelitynote}

%===================================================================
\section{Translation conventions}
%===================================================================

\begin{enumerate}
\item \textbf{Proof bodies omitted.}  Lean machine-checks all proofs.

\item \textbf{Definitional round trips.}  All four round-trip equalities for
  \lean{seqExt}/\lean{seqExtInv} and the two unitor/associator laws hold
  by \lean{rfl} after \lean{cases}, using only definitional reduction.

\item \textbf{Unicode notation.}  The Lean source uses the diamond operator
  $\triangleleft$ (\texttt{\\triangleleft}).  The listings use
  \texttt{seq} as an ASCII placeholder.

\item \textbf{Dependent type limitation.}  Haskell approximations drop
  the dependency of position type on shape value.
\end{enumerate}

\end{document}
